\documentclass[10pt]{article}
\usepackage[margin=20mm]{geometry}
\usepackage{amsmath,amssymb,mathtools,fontspec,unicode-math,xurl,hyperref}
\IfFontExistsTF{DejaVu Serif}{\setmainfont{DejaVu Serif}}{\setmainfont{Latin Modern Roman}}
\setmathfont{latinmodern-math.otf}
\setlength{\parindent}{0pt}\setlength{\parskip}{6pt}\sloppy
\title{Scientific review of M25-0400
Phase Transition of Parabolic Ginzburg–Landau Equation
with Potentials of High-Dimensional Wells}
\author{Initial-review recovery: exact-source mathematical audit}
\date{29 October 2026}
\begin{document}
\maketitle
Abstract: The reviewed source is Yuning Liu’s 38-page manuscript `arXiv:2207.12912v3`, dated 10 May 2023, associated with the frozen 2025 journal record, DOI `10.1002/cpa.22242`. Two major proof defects are established. The weighted normal spheres in Section 4 are used with Euclidean tube radii, invalidating the stated connection estimate and its use in the minimal-pair argument. The construction in Section 5 takes a constant normal extension of general \(H^{1/2}\) traces; an explicit admissible pair of \(H^1\) initial maps makes the constructed map fail to belong to \(H^1\). Eight further findings concern statements or calculations with bounded repair paths. Their scope is stated separately; in particular, a Hausdorff-convergence counterexample is at initial time, and a zero-gradient square-completion defect is an algebraic finding, not a claimed positive-time counterexample. No counterexample to the full main theorems is asserted. Substantial correction of the two major arguments is required for the reviewed version. Human adjudication remains pending. The current focused independent acceptance state is recorded in the accompanying provenance.\par
\section*{Identity, source boundary, and assessment}
The frozen annual identity is Yuning Liu, Phase Transition of Parabolic Ginzburg–Landau Equation with Potentials of High-Dimensional Wells, Communications on Pure and Applied Mathematics 78(6), 1199–1247 (2025), DOI 10.1002/cpa.22242 (https://doi.org/10.1002/cpa.22242). The author faculty page and publisher metadata support this identity. The actual mathematical source is the exact arXiv version below. Neither a change of annual identity nor equivalence between that manuscript and the journal version is inferred.\par
Item | Exact reviewed binding\par
Source | `papers/0400/0400.pdf`; 485,228 bytes; 38 pages\par
Version | arXiv:2207.12912v3 (https://arxiv.org/abs/2207.12912v3), 10 May 2023\par
Source SHA256 | `46dd24eb52d9897dd85570965841e4d477f93060284ca1d2836d86998879ad0b`\par
Manifest SHA256 | `161cf117f9d09ee93845505b348608c0e2c7ccc693e8c15b08ee58e726a74759`\par
Assessment | Two Major and eight Minor findings; full main result not refuted\par
Decision | `pending`; human reviews empty; current acceptance state in provenance\par
All page and equation references to the target below refer to this 38-page manuscript, not to the journal pagination. Index parsing remains separate from scientific reading. Before that original review, the available material contained source/version records and partial intake, but no substantive report or calculation for this paper. The historical false visual-inspection status remains preserved.\par
The two Major findings affect central claimed outputs: the proof of the minimal-pair boundary condition (1.19), and the general initial-data construction in Theorem 1.2. A finding need not refute every conclusion of a paper to be Major. The Minor findings are so classified because the downstream application has stronger hypotheses, the claim is unused, or a concrete local correction preserves the estimate.\par
\section*{Mathematical scope and dependency audit}
The equation is \[\partial_tu_\varepsilon=\Delta u_\varepsilon-\varepsilon^{-2}\partial F(u_\varepsilon),
 \qquad F(u)=f(\operatorname{dist}(u,m_+\sqcup m_-)^2),\] with smooth disjoint connected closed wells, \(0\in m_-\), smooth nearest-point projection in \(B_{2\delta_0}(m)\), and \[f(s)=s\quad(0\leq s\leq\delta_0^2),\qquad
 f(s)=2\delta_0^2\quad(s\geq2\delta_0^2).\] The domain is bounded with \(C^1\) boundary, the boundary datum is a smooth well-valued map into \(m_-\), and the prescribed smooth closed interface follows mean curvature flow while staying away from the outer boundary. Write \(D=\operatorname{dist}(m_+,m_-)\) and \[c_F=2\int_0^{D/2}\sqrt{2f(r^2)}\,\mathrm dr,\qquad
 \psi_\varepsilon=d_F(u_\varepsilon).\] The modulated energy is \[E_\varepsilon=\int_\Omega
 \left(\frac{\varepsilon}{2}|\nabla u_\varepsilon|^2+
       \frac{F(u_\varepsilon)}{\varepsilon}-\xi\cdot\nabla\psi_\varepsilon\right)\,\mathrm dx.\] The preparation assumption (1.15) bounds \(\varepsilon\|u_\varepsilon(0)\|_\infty+B(0)+E_\varepsilon(0)\) by \(C\varepsilon\). Theorem 1.1 asserts uniform energy and bulk-error control, phase convergence, subsequential harmonic heat-flow limits, and minimal-pair traces for every \(t\in(0,T)\). Theorem 1.2 quantifies over every pair of \(H^1\) well-valued maps whose interface traces are a minimal pair, and constructs bounded \(H^1\) approximants equal to those maps outside a fixed collar. These quantifiers matter in P1, P2, and P5.\par
Argument | Source pages | Audit result\par
SBV closure and compactness | 6, 20–22 | The stated jump-amplitude hypothesis is insufficient (P3); the later application also controls jump-set area.\par
Minimal connections and phase distance | 7–8 | The distinction between an infimum for arbitrary endpoints and attainment at minimal pairs is essential and consistent. The near-plus-well gradient sign needs correction (P4).\par
Interface geometry | 9–10 | The cutoff annulus invalidates one unrestricted normal derivative identity; its contracted version used later survives (P7).\par
Chain rule, projection, coercivity | 11–13 | The unconditional projected-gradient definition and the later classical-gradient convention differ (P4). The latter is explicitly printed and prevents a broader failure claim.\par
Energy and bulk error | 14–17, 32–37 | Square completion and dissipation were checked against primary relative-energy sources. P8 and P10 have bounded repairs.\par
Level sets and bulk limits | 18–22 | Perimeter/BV control does not identify the Hausdorff limit (P5). The normal-frame calculation needs its connection term (P6).\par
Minimal-pair traces | 22–26 | The weighted-radius connection estimate is invalid as written (P1); this is a central dependency of (1.19).\par
Initial preparation | 26–31 | General \(H^1\) data are not preserved by the collar construction (P2). The limiting bulk-energy coefficient also needs correction (P9).\par
References and versions | 37–38 | Key primary arguments were read with their hypotheses and normalizations; exact retrieved versions and limitations are retained.\par
\section*{P1: weighted and Euclidean radii in the connection estimate (Major)}
\subsection*{Location and claim}
The relevant definitions are (4.1)–(4.5), pp. 22–23; the affected statements are Lemmas 4.1 and 4.2, pp. 23–25. They feed directly into (4.16)–(4.20) and the contradiction proving (1.19), p. 26. The weighted semidistance and normal spheres are \[d_F^*(v,w)=\inf_\gamma\int\sqrt{2F(\gamma)}|\gamma'|\,\mathrm dt,\qquad
 N(p,r)=\{z:d_F^*(z,p)=r,\ P_{m_\pm}z=p\}.\] In (4.5), both normal spheres have weighted radius \(\delta_0\), and \(\kappa\) contains the additive correction \(2\delta_0-c_F\). Lemma 4.2 instead stops a curve when it exits the Euclidean tube \(B_{\delta_0}(m_\pm)\).\par
\subsection*{A counterexample to the stated lemmas}
Set \(a=\delta_0=1/100\) and take two unit circles in \(\mathbb R^2\), \[m_-=\{(1+\cos\theta,\sin\theta)\},\qquad
 m_+=\{(4+\cos\theta,\sin\theta)\}.\] They are compact, connected, smooth and disjoint; \(0\in m_-\), \(D=1\), and the \(2a\)-tubes have smooth nearest-point projections. The closest pair is \(p_-=(2,0)\), \(p_+=(3,0)\).\par
Choose a monotone smooth \(f\) satisfying the prescribed endpoint conditions, with \(a^2\leq f\leq2a^2\) in the intermediate interval. For completeness, such a \(C^\infty\) choice can be made explicitly. For \(0<z<1\), put \[s(z)=\frac{e^{-1/z}}{e^{-1/z}+e^{-1/(1-z)}},\quad
 b(z)=e^{-1/[z(1-z)]},\quad B=\int_0^1b(z)\,\mathrm dz,
 \quad w(z)=1-s(z)+\frac{b(z)}{2B}.\] Extend \(s\) by \(0,1\) and \(b\) by zero at the endpoints. Symmetry gives \(\int_0^1s=1/2\), hence \(\int_0^1w=1\); \(w\geq0\), \(w(0)=1\), \(w(1)=0\), and all derivatives needed for gluing are flat. Define \[f(t)=a^2\left[1+\int_0^{(t-a^2)/a^2}w(z)\,\mathrm dz\right]
       \quad(a^2<t<2a^2)\] and use \(f(t)=t\) below and \(f(t)=2a^2\) above. Thus the example satisfies the potential hypotheses, including smoothness of \(F\) across the parts where the distance function has a cut locus, because \(f\) is already constant there.\par
Equation reference: eq:cfbounds. Direct bounds on the defining integral give \[2a-4\sqrt2\,a^2\leq c_F
 \leq 2a-(4-\sqrt2)a^2.\] In particular, \[
 0.019434\leq c_F\leq0.0197415,\qquad a<c_F<2a.\] The displayed decimal bounds are certified by the rational enclosure \(1.414<\sqrt2<1.415\), not by floating-point quadrature.\par
On the closest segment, write \(z(q)=(2+q,0)\), \(0\leq q\leq1\), and \[S(q)=\int_0^q\sqrt{2f(\min\{r,1-r\}^2)}\,\mathrm dr.\] Then \(S(0)=0\), \(S(1)=c_F\), and \(d_F^*(z(q),p_-)=S(q)\), \(d_F^*(z(q),p_+)=c_F-S(q)\). Here is an independent justification of the lower bounds implicit in these equalities. Put \(h(x)=\min\{1,\max\{0,x_1-2\}\}\). For \(2\leq x_1\leq3\), the distance to the two circles is at least \(\min\{h(x),1-h(x)\}\). Monotonicity of the chosen \(f\) therefore gives \(|\nabla(S\circ h)|\leq\sqrt{2F}\) almost everywhere. Integrating this calibration along any path gives the lower bound; the straight segment attains it. This also excludes shortcuts along the zero-cost wells.\par
Equation reference: eq:kappapos. Because \(a<c_F\), the spheres \(N(p_\pm,a)\) are nonempty, with explicit points on this segment. Their distance is finite and nonnegative. For every sufficiently small \(\rho\geq0\), the definition (4.5) therefore implies \[
 \kappa(p_+,p_-,\rho)\geq2a-c_F
 \geq \frac{517}{2000000}>0.\] At \(\rho=0\) this contradicts the zero-value conclusion of Lemma 4.1 for the minimal pair \((p_+,p_-)\).\par
It also contradicts Lemma 4.2 itself. Let \(Q:\mathbb R\to(0,1)\) be the increasing optimal profile \[Q'=\sqrt{2f(\min\{Q,1-Q\}^2)},\qquad Q(-\infty)=0,\quad Q(+\infty)=1,\] and set \(\gamma(\tau)=z(Q(\tau/\varepsilon))\) on a fixed interval \([-h,h]\). For small \(\varepsilon\), \(\gamma(\pm h)\in B_a(m_\pm)\), so (4.7) holds. Along this curve \(d_F(\gamma)=S(Q)\), and pointwise \[\frac12|\gamma'|^2+\frac{F(\gamma)}{\varepsilon^2}
       -\frac{(d_F\circ\gamma)'}{\varepsilon}=0.\] For every fixed \(0<\rho<\rho_0\), the claimed lower bound (4.8) is strictly positive by [eq:kappapos] (\#eq:kappapos), whereas its left side is zero. This is an exact counterexample, with no numerical solution of the author’s PDE or code.\par
\subsection*{Scope of a possible smallness interpretation}
The proof of Lemma 3.4 mentions \(\delta_0\ll c_F\) on p. 18. The unit-gap example does not satisfy that extra informal relation. It does satisfy the explicitly stated hypotheses (1.3)–(1.6). Because the same \(\delta_0\) fixes the height of the plateau in \(f\), shrinking \(\delta_0\) in this example does not make \(\delta_0/c_F\) tend to zero. Thus this relation cannot silently be obtained from the written assumptions by making the tube thinner.\par
There is also a separate defect even under an arbitrarily strong relation \(\delta_0\ll c_F\). Move the second circle’s center to \((L+3,0)\), so that \(D=L\) can be as large as desired. Then \(c_F\sim2aL\), but an actual Euclidean tube exit at distance \(a\) has weighted distance \[r_a=\int_0^a\sqrt2\,r\,\mathrm dr=\frac{a^2}{\sqrt2}\neq a.\] After correcting the exit-time notation, the optimal straight profile has \(P_m\gamma(t_+)=P_m\gamma(h)=p_+\) and \(\gamma(t_+)\notin N(\overline B_\rho(p_+)\cap m_+,a)\). Case (1) on p. 25 would nevertheless force \(|P_m\gamma(t_+)-P_m\gamma(h)|\geq C_1\rho>0\). This is false even when \(a/c_F\) is arbitrarily small.\par
Two additional local errors are part of this same proof failure. The printed infimum defining \(t_-\) on p. 24 equals \(-h\) rather than the intended first exit; the printed supremum defining \(t_+\) equals \(h\) rather than the last entrance. The infimum and supremum must be interchanged. Also, the upper bound ending \(d_F^*(N_+,N_-)\leq c_F-2\delta_0+\beta\) on p. 24 does not imply the asserted equality \(\kappa=\beta>0\). The first equality for \(I_0\) in case (3), p. 25, is in general a lower bound obtained by Cauchy–Schwarz.\par
\subsection*{Impact and repair boundary}
The estimate used to force minimal-pair traces is not established. This is a Major issue because (1.19) is a principal advertised conclusion. The counterexample concerns the auxiliary connection lemma, not an actual limiting pair violating Theorem 1.1.\par
A consistent reconstruction must use a geometric tube radius \(r\) and the distinct weighted radius \(r_F=\int_0^r\sqrt{2f(s^2)}\,\mathrm ds\). The normal spheres and additive correction must then use \(r_F\) and \(2r_F\), while the exit times use the Euclidean tube of radius \(r\). The strict gap for nonminimal endpoint pairs, its compactness and continuity properties, and the resulting curve estimate still require a complete proof. Merely replacing one occurrence of \(\delta_0\) does not verify this repair. The precise minimal-connection equality case in Lin–Pan–Wang [LPW], Section 2, is relevant, but does not by itself validate the printed \(\kappa\) argument.\par
\section*{P2: the collar construction does not preserve general \(H^1\) data (Major)}
\subsection*{Location}
On p. 26, (5.1) calls \(\Psi_\delta^\pm\) global \(C^1\) diffeomorphisms, while requiring \(\Psi_\delta^\pm=P_{\Sigma_0}\) on the entire open collar \(B_\delta(\Sigma_0)\). On p. 27, (5.2) asserts that \(u_0^\pm=u_{\pm}^{\rm in}\circ\Psi_\delta^\pm\) belongs to \(H^1\), and (5.10) uses these extensions to define \(u_\varepsilon^{\rm in}\).\par
Already at the smooth geometric level, \(DP_{\Sigma_0}\) has rank \(d-1\) in the collar. Such a map is not a diffeomorphism between \(d\)-dimensional domains. The more substantial defect remains even if the map is renamed a collapsing extension.\par
\subsection*{Admissible trace counterexample}
Take \(\Omega=B_3(0)\subset\mathbb R^2\) and \(\Sigma_0=S^1\), with \(\Omega_0^+=B_1(0)\) and \(\Omega_0^-=B_3(0)\setminus\overline B_1(0)\). A short smooth mean curvature evolution exists, and the separation from the outer boundary meets the geometric hypothesis for small \(\delta_0\). In target space \(\mathbb R^3\), take the parallel circles \[m_-=\{(1-\cos\theta,\sin\theta,0)\},\qquad
 m_+=m_-+(0,0,D),\quad D>0.\] They satisfy all well hypotheses, including \(0\in m_-\), and pairs at the same angle have exactly the minimal distance \(D\).\par
Let \[a(\theta)=c\sum_{j=1}^{\infty}2^{-3j/4}\sin(2^j\theta),
 \qquad c=\frac1{100}.\] The series converges absolutely and uniformly and has \(|a|<\pi/8\). Its Sobolev seminorms satisfy \[\sum_j 2^j\,2^{-3j/2}<\infty,\qquad
 \sum_j 2^{2j}\,2^{-3j/2}=\infty.\] Thus \(a\in H^{1/2}(S^1)\cap C(S^1)\) but \(a\notin H^1(S^1)\). This is a divergence proof, not an inference from finitely many sampled Fourier modes.\par
Let \(a_+\) be the harmonic extension to the disk, and let \(a_-\) be the harmonic extension to the annulus with inner boundary value \(a\) and outer boundary value zero. The Fourier energy formulas, or the \(H^{1/2}\) trace theorem, give \(a_\pm\in H^1\). The maximum principle keeps their values in \((-\pi/8,\pi/8)\). Define \[u_-^{\rm in}=(1-\cos a_-,\sin a_-,0),\qquad
 u_+^{\rm in}=(1-\cos a_+,\sin a_+,D).\] These are bounded \(H^1\) well-valued maps. The outer boundary value is the smooth constant zero, and their interface traces have distance exactly \(D\). In particular, they satisfy the full minimal-pair assumption (1.23), not merely separate phase membership.\par
On the collar, (5.1)–(5.2) give the first two coordinates \[(1-\cos a(\theta),\sin a(\theta)),\] constant in the normal coordinate. This is not an \(H^1\) map on any positive-width annulus. If it were, its second component would be \(H^1\) there, and composition with the Lipschitz function \(\arcsin\) on its range would give the constant normal extension of \(a\) in \(H^1\). Integration in the normal variable would imply \(a\in H^1(S^1)\), a contradiction.\par
The failure persists in the actual final construction (5.10). The two extended maps have identical first two coordinates and differ only in the third coordinate. Therefore interpolation by \(S_\varepsilon\) changes only that third coordinate. The non-\(H^1\) first two coordinates remain unchanged for every \(\varepsilon\). The constructed \(u_\varepsilon^{\rm in}\) is consequently not in \(H^1\).\par
\subsection*{Impact and repair boundary}
This disproves the claimed construction for admissible inputs of Theorem 1.2. It does not prove that no other preparation sequence exists. A valid proof needs a Sobolev-compatible gluing construction that preserves the minimal-pair constraint and the \(O(\varepsilon)\) modulated-energy bound, possibly with smoothing scales depending on \(\varepsilon\). Alternatively, stronger tangential trace regularity would change the theorem’s stated scope. Renaming \(\Psi_\delta^\pm\), or citing a trace theorem that supplies only \(H^{1/2}\) traces, does not repair the argument. The initial construction in the separately checked Laux–Liu preprint [LL] starts with a global \(H^1\) map and does not supply this missing extension for arbitrary pairs of bulk maps.\par
\section*{Further confirmed findings with bounded repair paths}
\subsection*{P3: insufficient SBV jump hypothesis (Minor)}
Location. Propositions 2.2 and 2.3, p. 6, especially (2.2); application in Theorem 3.6, p. 22.\par
Counterexample. On \((0,1)\), let \(\mu_j\) put mass \(2^{-j}\) at the midpoint of each of the \(2^j\) surviving intervals in the \(j\)th middle-third Cantor construction, and let \(v_j(x)=\mu_j((0,x])\). Each \(v_j\) is a step function in \(\operatorname{SBV}\), \(0\leq v_j\leq1\), \(\nabla^a v_j=0\), and \[\int_{J_{v_j}}|v_j^+-v_j^-|\,\mathrm d\mathcal H^0=1.\] Choose \(\varphi(t)=t^2\), which satisfies all conditions in Proposition 2.2. The distribution functions converge uniformly to the continuous Cantor function \(v\), and their derivative measures converge weakly to Cantor measure. Thus \(v_j\to v\) weakly star in \(BV\), but \(D^cv\neq0\) and \(v\notin\operatorname{SBV}\). Every assumption in both stated propositions is met.\par
Impact and correction. A linear jump-amplitude cost does not exclude a Cantor part. An unweighted uniform bound on jump-set measure, together with the other displayed bounds, is sufficient in the actual application. More generally, a jump cost \(\vartheta\) with \(\vartheta(s)/s\to\infty\) as \(s\downarrow0\) supplies the missing small-jump control; see the precise primary theorem [AM], Theorem 1.4 and Proposition 2.5. Theorem 3.6 already invokes the area bound (3.28a), the gradient bound (3.29b), and bounded values. Replacing the quoted compactness theorem by its correct form is therefore a bounded repair, which is why this is Minor despite the false preliminary statements.\par
\subsection*{P4: incompatible gradient conventions and a well sign (Minor)}
Location. The plus-well case of (2.21), p. 8; the unconditional definition (2.35) and the classical-coincidence assertion in Lemma 2.7, p. 11; the convention immediately below (2.40), p. 12; and (2.41).\par
At a differentiability point, (2.35) projects the classical gradient of \(d_F\) onto \[V=\operatorname{span}\{\partial_tu_\varepsilon,\partial_1u_\varepsilon,
                  \ldots,\partial_du_\varepsilon\}.\] Such a projection need not equal the full classical gradient. This is not merely a possible choice of arbitrary jets: the following exact smooth solution of (1.7) demonstrates it.\par
Use \(\Omega=(0,\pi)\) and two unit circles in the first two coordinates of \(\mathbb R^5\), centered at \((1,0)\) and \((5,0)\); their last three coordinates are zero. Set \(g=0\), \(\delta_0=0.01\), \(\varepsilon=0.1\), \(K=100\), \(t_*=K^{-2}\), and \(A=\delta_0 e^{-10}/100\). For \(i=1,2,3\), put \[A_i=Ae^{(2/\varepsilon^2+(iK)^2)t_*},\qquad
 u_\varepsilon(x,t)=
 \left(0,0,\left(A_i e^{-(2/\varepsilon^2+(iK)^2)t}\sin(iKx)\right)_{i=1}^3\right).\] The initial amplitude norm is less than \(3.754\times10^{-5}\), well inside the \(0.01\) tube, and decreases in time. The nearest point on \(m_-\) is zero, so \(F(u_\varepsilon)=|u_\varepsilon|^2\); hence this formula solves the actual linearized equation \(\partial_tu_\varepsilon=\partial_{xx}u_\varepsilon-2\varepsilon^{-2}u_\varepsilon\), which is exactly (1.7) throughout its range. The boundary value is zero. No author code or numerical PDE solver is used.\par
At \(x_*=\pi/(4K)\), \(t=t_*\), set \(b=2/(\varepsilon^2K^2)=1/50\) and \[U=(1/\sqrt2,1,1/\sqrt2),\quad
 V_1=(1/\sqrt2,0,-3/\sqrt2),\quad
 W=((b+1)/\sqrt2,b+4,(b+9)/\sqrt2).\] Then \(u=AU\), \(u_x=AKV_1\) and \(u_t=-AK^2W\) in the last three coordinates. Writing \(D_b=9b^2+64b+116\), exact Gram-matrix calculation yields \[|\operatorname{proj}_{\operatorname{span}(V_1,W)}U|^2
       =2-\frac4{D_b}<2.\] Near \(m_-\), \(d_F(u)=|u|^2/\sqrt2\) and \(\nabla d_F=\sqrt2u\). Thus the literal (2.35) gradient \(g_{\rm pr}\) has \[|g_{\rm pr}|^2=4A^2-\frac{8A^2}{D_b}
     =4A^2\frac{288209}{293209},\] whereas \(|\nabla d_F|^2=4A^2\). The radial projection prescribed in the well tube has \(|\Pi_u u_x|^2=A^2K^2/2\), and therefore \[|\psi_x|^2=2A^4K^2,\qquad
 |\Pi_u u_x|^2|g_{\rm pr}|^2
   =A^4K^2\left(2-\frac4{D_b}\right).\] This contradicts (2.41) if the unconditional (2.35) convention is used there. It also directly disproves the asserted classical coincidence of that projected definition.\par
The scope qualification is decisive: the actual text immediately below (2.40) says to use the generalized differential when \(d_F\) is not classically differentiable. With the classical gradient at differentiability points, this example no longer contradicts (2.41). That explicit later convention was checked in the source image and is a local repair. The example is not advanced as a counterexample to the well-prepared limit theorem or to its integrated energy estimate.\par
Finally, differentiating (2.17) near the plus well gives \[\nabla d_F(u)=-\sqrt2\,(u-P_{m_+}u),\] while near the minus well the sign is positive. The common positive sign printed in (2.21) needs correction. Projection onto the resulting line is unchanged by this sign. State one consistent piecewise gradient convention, and correct the sign and the classical-coincidence claim.\par
\subsection*{P5: initial preparation permits remote small droplets (Minor)}
Location. Proposition 3.5, (3.28c), pp. 19–20, and the appeal to Blaschke’s theorem in its proof.\par
Admissible initial-time counterexample. Let \(d=3\), \(\Omega=B_5(0)\), and let the prescribed interface start at \(\Sigma_0=\partial B_2(0)\) and follow \(\Sigma_t=\partial B_{\sqrt{4-4t}}(0)\) for a short time. Use the compact two-circle wells from P1 and the closest segment between \(p_-\) and \(p_+\), with a smooth optimal transition profile. The outside boundary value is \(p_-\), which is an admissible smooth map into \(m_-\). A truncated optimal layer across \(\Sigma_0\) gives \(E_\varepsilon(0)=O(\varepsilon^2)\) and \(B(0)=O(\varepsilon^2)\) when the bulk phases are constant.\par
Add a second, plus-phase ball centered at \(x_0=(3,0,0)\), of radius \(r_\varepsilon=\sqrt\varepsilon\), inside the minus phase. Use the same smooth profile across this sphere, with thickness of order \(\varepsilon\), and truncate its tails while staying away from the principal interface and the outer boundary. The calibration \(\xi\) vanishes near the small ball. Its additional energy and bulk error satisfy \[E_{\varepsilon,\mathrm{drop}}
  =O(r_\varepsilon^2+\varepsilon r_\varepsilon+\varepsilon^2)=O(\varepsilon),\qquad
 B_{\mathrm{drop}}=O(r_\varepsilon^3+\varepsilon r_\varepsilon^2)=O(\varepsilon^{3/2}).\] Exponential tail errors do not change these bounds. The values remain on the bounded closest segment. Consequently the full preparation condition (1.15) holds with a uniform constant. These are smooth compatible initial data for (1.7).\par
For every sufficiently small \(b_\varepsilon\in[\varepsilon^{1/6},2\varepsilon^{1/6}]\), the level set \(\{\psi_\varepsilon=c_F-b_\varepsilon\}\) has a component around \(x_0\). Its displacement from the small sphere is \(O(\varepsilon|\log\varepsilon|)=o(r_\varepsilon)\). Points on this component remain a fixed positive distance from \(\Sigma_0\). No subsequence can satisfy the asserted Hausdorff convergence at \(t=0\). At the same time, its extra perimeter is \(O(\varepsilon)\) and its enclosed volume is \(O(\varepsilon^{3/2})\), fully consistent with the perimeter and BV conclusions.\par
Scope and repair. This is an initial-time counterexample under the stated uniform preparation assumptions. It does not show a remote droplet surviving to a fixed positive time: small droplets may disappear rapidly. Blaschke compactness only supplies some Hausdorff-convergent subsequence of compact sets; it does not identify the limit from perimeter and \(L^1\) convergence. Theorem 3.6 uses (3.28a)–(3.28b), and the minimal-pair proof uses slice convergence, traces and Egorov’s theorem, rather than (3.28c). Removing that unused claim, or replacing it by a separately proved positive-time clearing statement, is a bounded repair.\par
\subsection*{P6: normal-frame derivatives in higher codimension (Minor)}
Location. The normal-coordinate calculation (3.34)–(3.37), pp. 20–21, used for (3.29b) and again in (4.10).\par
Let \(q(s)=(\cos s,\sin s,0)\) and choose the smooth normal frame \[\nu_1(s)=\cos s\,q(s)+\sin s\,e_3,\qquad
 \nu_2(s)=-\sin s\,q(s)+\cos s\,e_3.\] For fixed small \(r>0\), set \(u(s)=q(s)+r\nu_2(s)\). The nearest point is \(q(s)\) and the normal coordinate \(d_1(u(s))=\langle u(s)-q(s),\nu_1(s)\rangle\) vanishes identically. But at \(s=0\), \[\frac{\,\mathrm d}{\,\mathrm ds}d_1(u(s))=0,\qquad
 \nu_1(0)\cdot u'(0)=-r\neq0.\] Thus the asserted gradient formula for individual normal coordinates, and the claim that all derivatives of a normal frame are tangential, fail in general codimension. Translating the circle makes \(0\in m_-\); a second distant circle supplies the other well. The example is within a smooth tubular neighborhood.\par
The required estimate can be proved without this false frame identity. If \(\nu_{\rm rad}=(u-P_mu)/|u-P_mu|\), then \(DP_m(u)\nu_{\rm rad}=0\), and smooth tube geometry gives \(\|DP_m(u)\|\leq1+C\delta_0\). Hence \[|\partial_iP_m(u)|
 \leq(1+C\delta_0)
     |\partial_i u-\Pi_{\rm rad}\partial_i u|.\] This supplies the bounded projection estimate used later. Correcting the normal connection, or using this coordinate-free proof, avoids a loss in the essential energy bound.\par
\subsection*{P7: curvature extension through its cutoff annulus (Minor)}
Location. The first identity in (2.31), p. 10, compared with (2.30).\par
For a circle or sphere of nonzero mean curvature, the extended curvature field has the form \(H(x)=\kappa(P_{\Sigma_t}x)\eta_0(x)n(x)\). Choose an admissible radial cutoff, equal to one for \(|d_\Sigma|\leq\delta_0\) and zero for \(|d_\Sigma|\geq2\delta_0\). On a normal ray, \[(n\cdot\nabla)H=\kappa\,\eta_0'(d_\Sigma)n.\] The derivative is nonzero somewhere in the cutoff annulus. For the flat switch \(s\) from P1, choose a smooth \(\chi(r)\) equal to one for \(r\leq7a/4\) and zero for \(r\geq15a/8\). Set \(\eta_0(r)=[1-s((r-a)/a)]\chi(r)\) for \(a<r<2a\), extend by one for \(0\leq r\leq a\), then evenly. Its support lies strictly inside \((-2a,2a)\), as required by (2.29), while \(\eta_0'(3a/2)=-2/a\) is unchanged. Thus the unrestricted normal derivative identity cannot hold throughout \(B_{2\delta_0}(\Sigma_t)\).\par
On the smaller tube where \(\eta_0=1\), the normal derivative does vanish. Since \(\xi\) is supported in that smaller tube, \((\xi\cdot\nabla)H=0\) remains valid everywhere it is needed. The appendix uses this contracted relation. Restrict the first identity’s domain and retain the contracted identity.\par
\subsection*{P8: derivative of the negative part (Minor)}
Location. The displayed calculation (3.13), p. 16, with the definitions immediately preceding it.\par
The convention is the positive-valued negative part \(z^-=\max\{-z,0\}\). Away from the interface and where \(0<\psi_\varepsilon<c_F\), the function \[w_\varepsilon=c_F\chi-c_F+2(\psi_\varepsilon-c_F)^-
        =c_F\chi+c_F-2\psi_\varepsilon\] satisfies \(\partial_iw_\varepsilon=-2\partial_i\psi_\varepsilon\). The printed plus sign is incorrect. A nonconstant monotone optimal profile at a point strictly inside one phase has \(\partial_i\psi_\varepsilon\neq0\) and \(\eta(d_\Sigma)\neq0\), so the discrepancy is realized by smooth admissible local data, rather than only at the nondifferentiable threshold.\par
The subsequent estimates use absolute values and Young’s inequality. Correcting this sign leaves those bounds available. Also, the positive-part auxiliary \((\psi_\varepsilon-c_F)^+\) is identically zero for the bounded phase distance used here; introducing it is redundant but is not itself a mathematical defect.\par
\subsection*{P9: the bulk-energy coefficient in (5.40) (Minor)}
Location. Equations (5.39)–(5.40), p. 31.\par
Use smooth well-valued bulk maps that are constant near \(\Sigma_0\), satisfy the same-angle minimal-pair condition for the parallel circles in P2, and are nonconstant only on a compact subset of one bulk region outside \(B_{2\delta}(\Sigma_0)\). For example, take a nonconstant smooth compactly supported angle function in the plus bulk and the zero angle elsewhere. These data avoid the trace-regularity problem in P2 altogether.\par
Let \(R=\sum_\pm\int_{\Omega_0^\pm}|\nabla u_0^\pm|^2>0\). The construction is the exact original bulk map on the support of those gradients, where \(F=0\) and \(\nabla\psi=0\). Its contribution to \(E_\varepsilon/\varepsilon\) is exactly \(R/2\). The optimal normal layer contributes \(O(\varepsilon)\) after this division, and the truncated tails tend to zero exponentially. Thus \[\lim_{\varepsilon\downarrow0}\frac{E_\varepsilon[u_\varepsilon^{\rm in}|\Sigma_0]}{\varepsilon}
    =\frac12\sum_\pm\int_{\Omega_0^\pm}
                  |\nabla u_0^\pm|^2
    =\frac R2,\] whereas (5.40) prints the sum without the factor \(1/2\). This is also the coefficient carried by the preceding formula (5.39). Restoring it does not alter the \(O(\varepsilon)\) preparation estimate.\par
\subsection*{P10: zero spatial gradient in the final square completion (Minor)}
Location. The zero convention in (2.39a), p. 11, and the final square-completion step on p. 35 leading to (A.2a), p. 33.\par
Write \(g=\partial d_F(u_\varepsilon)\) and \(c=\operatorname{div}\xi\). Immediately before its last square, the proof contains \[c^2|g|^2+2\varepsilon c\,\nabla\psi_\varepsilon\cdot H
             +\varepsilon^2|\Pi_{u_\varepsilon}\nabla u_\varepsilon|^2|H|^2.\] It replaces this by \(|c|g|n_\varepsilon+\varepsilon|\Pi_{u_\varepsilon}\nabla u_\varepsilon|H|^2\). Where \(\nabla\psi_\varepsilon=0\), the printed convention is \(n_\varepsilon=0\). If \(g\neq0\) and \(\Pi_{u_\varepsilon}\nabla u_\varepsilon=0\), the first expression is \(c^2|g|^2\) and the proposed square is zero. The chain rule does not remove that discrepancy.\par
These pointwise data are compatible with smooth prepared initial data. Start with a truncated optimal layer and choose a small fixed open patch \(U\) in the minus side of the calibration collar, away from the actual layer, where \(\operatorname{div}\xi\neq0\). For the parallel circles in P2, take the closest segment to be vertical and add the smooth perturbation \(\varepsilon\zeta e_3\) to the constant minus phase, where \(\zeta=1\) on \(U\) and is supported in that side of the collar. For small \(\varepsilon\), this adds \(O(\varepsilon)\) to the energy and \(O(\varepsilon^2)\) to the bulk error, and changes no outer boundary value. On \(U\) at \(t=0\), \[u_\varepsilon=\varepsilon e_3,\quad \nabla u_\varepsilon=0,\quad
 F=\varepsilon^2,\quad g=\sqrt2\varepsilon e_3,\quad
 \partial_tu_\varepsilon=-2\varepsilon^{-1}e_3.\] The omitted term after division by \(2\varepsilon\) is \(\varepsilon(\operatorname{div}\xi)^2>0\).\par
Precise scope. This establishes a defect in the unconditional algebraic completion, with an admissible smooth initial plateau. It is not claimed to construct a fixed positive-time plateau or to refute an almost-everywhere interior-time version of Lemma A.2. Such a stronger assertion would require additional analysis.\par
A bounded correction is explicit. Choose an arbitrary unit normal where \(\nabla\psi_\varepsilon=0\), as in Fischer–Laux–Simon [FLS], (11a), or retain the omitted nonnegative remainder. On that zero-gradient set, \[\frac{c^2|g|^2}{2\varepsilon}
 \leq \|\operatorname{div}\xi\|_\infty^2\,\frac{F(u_\varepsilon)}{\varepsilon}
 \leq C\left(\frac{\varepsilon}{2}|\nabla u_\varepsilon|^2+
                    \frac{F(u_\varepsilon)}{\varepsilon}
                    -\xi\cdot\nabla\psi_\varepsilon\right).\] Thus the remainder is controlled by \(E_\varepsilon\) and does not by itself refute Proposition 2.9. The three-square estimate on p. 36 should also use a generic constant (for example the usual factor \(3\)), rather than coefficient one; that constant is absorbed in the same estimate.\par
\section*{Checks that do not support additional confirmed failures}
\subsection*{Infimum versus attainment of a connection}
The value \(C_F(p_+,p_-)=c_F\) for arbitrary well endpoints does not imply that a least-energy heteroclinic joins every such pair. One can move along either connected well at zero potential cost, taking arbitrarily long time to make its kinetic cost arbitrarily small, and use a closest-pair transition in between. This proves the common infimum. Attainment at the least energy forces the straight minimal-pair orbit. The distinction in the checked Lin–Pan–Wang equality-case argument is compatible with the target’s preliminary discussion. It does not cure P1’s incorrect radii.\par
\subsection*{The corrected dissipation term in the appendix}
The target explicitly retains \[\frac1{2\varepsilon}\int_\Omega
       \bigl(\varepsilon^2|\partial_tu_\varepsilon|^2-|H_\varepsilon|^2\bigr)\,\mathrm dx\] in Lemma A.2. In the exact Laux–Liu arXiv v2 checked here, the corresponding displayed Lemma 4.4 omits this term, while its square-completion calculation (4.35) contains it. The target notices this discrepancy. The reference’s statement therefore cannot be used to accuse the target of omitting that same term. The chain-rule convention in P4 and the zero-normal convention in P10 are separate questions.\par
\subsection*{Time quantifiers for the limiting traces}
A spatial \(H^1\) conclusion initially obtained for almost every time need not prevent a properly chosen representative from having minimal-pair traces at every time. Away from the smooth moving interface, the strong \(C_tL^2_{\rm loc}\) convergence and bounded values give global \(L^2\) continuity after pulling back the moving domains: the contribution of a collar of width \(\delta\) is bounded by its volume, uniformly in time. The uniform \(H^1\) estimate then extends to every time by weak lower semicontinuity.\par
On a fixed smooth reference domain the trace interpolation estimate \[\|\operatorname{Tr}v\|_{L^2(\partial\Omega)}^2
 \leq C\|v\|_{L^2(\Omega)}\|v\|_{H^1(\Omega)}\] gives continuity of the traces in \(L^2\). The minimal-pair constraint is closed under that convergence. Consequently an almost-everywhere-time result can be extended once the representative and these steps are supplied. This is a repair of presentation and quantifiers, not a replacement for the missing connection estimate in P1.\par
The level choices \(b_k^\pm\) in Proposition 3.5 should likewise make their possible time dependence explicit. A selection made by Fubini at each fixed time does not itself prove the existence of one threshold working at all times. No full counterexample to that stronger uniform selection was established in this review, so it is not counted as a confirmed additional defect.\par
\subsection*{The level-perimeter rate can be recovered by clipping}
The Chebyshev step on p. 19 does not display the dependence on the small level parameter. The desired level-selection bound can be obtained without relying on that omission. Let \(0<2\delta<c_F\) and define \[T_\delta(z)=\min\{\max\{z-c_F+2\delta,0\},\delta\}.\] With \(A_t=\mathcal H^{d-1}(\Sigma_t)\), the calibration and the \(L^1\) phase-error estimate give \[\begin{aligned}
 \int_{\{c_F-2\delta<\psi<c_F-\delta\}}|\nabla\psi|
 &\leq E_\varepsilon+\int_\Omega\xi\cdot\nabla T_\delta(\psi)\\
 &\leq C\varepsilon+\delta A_t+C\|\psi-c_F\mathbf 1_{\Omega_t^+}\|_{L^1}\\
 &\leq \delta A_t+C\sqrt\varepsilon.\end{aligned}\] Here integration by parts uses \(T_\delta(c_F\mathbf 1_{\Omega_t^+})=\delta\mathbf 1_{\Omega_t^+}\) and the Lipschitz constant one of \(T_\delta\). The coarea formula gives the needed averaged upper bound. For an individual regular level \(s\) in the same interval, \[P(\{\psi>s\};\Omega)
 \geq-\int_{\{\psi>s\}}\operatorname{div}\xi
 \geq A_t-C\frac{\sqrt\varepsilon}{\delta}.\] Combining the two bounds selects a level with the claimed order \(\sqrt\varepsilon/\delta\). Setting \(\delta=\varepsilon^{1/6}\) gives order \(\varepsilon^{1/3}\). This bounded derivation supports the perimeter route; it supplies no Hausdorff control, and hence does not repair P5.\par
\subsection*{Passage to harmonic heat flow}
Weak \(H^1\) convergence alone would not justify passage through the quadratic second fundamental form. The primary Chen–Struwe argument [CS], Sections II–IV, instead uses a penalized Bochner estimate, a localized monotonicity formula, small-energy regularity, normality of the limiting penalization force off a concentration set, and a cutoff removal argument. The target’s local energy and time-dissipation estimates (3.2), together with its bounded values, supply the type of local inputs needed away from the moving interface. Applying the argument requires localization to fixed space-time cylinders separated from that interface; the whole unscaled energy of the target is not uniformly bounded across the transition layer. This dependency was checked at that local level. No counterexample to the harmonic-flow passage was established. A full independent redevelopment of the entire partial-regularity theory is not claimed.\par
\section*{What a corrected argument would need to establish}
For P1, a replacement quantitative connection lemma must use compatible radii, prove a strict gap for nonminimal endpoint pairs, and give the uniform compactness needed on the positive-measure trace set used in (4.19). It must then be carried through the slicing contradiction. For P2, a new preparation construction must work for the full \(H^1\) class in (1.23), retain the prescribed maps outside \(B_{2\delta}(\Sigma_0)\), and control both \(E_\varepsilon\) and \(B\) by \(C\varepsilon\). Neither task is completed by the local corrections listed for the Minor findings.\par
The smaller corrections are more sharply delimited: use a valid SBV compactness statement with the available jump-area bound; make the gradient convention consistent; remove or separately prove the Hausdorff assertion; use the coordinate-free projection bound; restrict the curvature identity to its valid region; correct the negative-part sign and bulk coefficient; and make the zero-gradient square completion exact, or estimate its explicit remainder. None requires a claim that the authors have already adopted a correction.\par
\section*{Evidence, computation, and recovery record}
Its substantive report and evidence are preserved. Recovery-0027 subsequently completed the compilation, PDF parsing, report-page inspection, and delivery metadata without repeating the full scientific review. They did not perform science, calculations, compilation, report authorship, or image inspection. The controller’s role remains scheduling and byte collection. Those actions completed the initial delivery; the current focused acceptance is recorded separately in provenance.\par
The exact source text extraction is bound by SHA256 `9623ca77c72884fdf3b618b5ee33f7a6a9f8f3bf1861768db8a51a11ba6feb4a`. Complete primary texts were read for Chen–Struwe, Alberti–Mantegazza, Fischer–Laux–Simon arXiv v2, and Laux–Liu arXiv v2. For the 56-page Lin–Pan–Wang preprint, the declared reading scope is the identity and hypotheses, complete Section 2, and complete Appendix A. Other sections of that reference are not claimed to have been fully read.\par
The independent script uses exact rational arithmetic for the certified radius bounds, Gram projection, Cantor approximants, and coefficient checks. Decimal evaluations only supplement the analytic proofs. Ten issue-specific calculation records passed their assertions; this is not described as automated verification of the full paper. Their result file has SHA256 `6cea2db37ce37c2c9306e62014dfc6e1d1b1a74cc7499d497ff61718bfbfc132`. No author code was executed.\par
Retained scientific evidence is under `verification/0400/recovery\_e62ac067/`, including `independent\_checks/calculations.json`, `actual\_source\_views.json`, `actual\_primary\_views.json`, and `primary\_read\_coverage.json`. The companion coverage, findings, provenance and support manifest bind the report’s claims, source locations and actual supporting paths, byte counts and SHA256 values. Current build receipts bind the exact TeX, source, retained Tectonic binary, dependency payloads, logs and delivered PDF. Render receipts and actual page observations remain separate. Unchanged source and reference inspections are reused only through verified image and source hashes, preserving the original performing role. Original execution paths, the failed build and failed acquisition records remain immutable historical evidence.\par
It refines the P7 cutoff to have support strictly inside \(B_{2a}\); the local derivative and Minor rating are unchanged. Its current rebuild, page checks, and detailed acceptance provenance are bound under `verification/0400/acceptance\_1c9bcad6/`.\par
A bounded public search of the exact title, identifier, publisher and author pages did not identify a separate applicable supplement or code package. This is not a proof that none exists, and missing code is not a mathematical finding. The publisher download for the Laux–Liu reference returned a non-PDF payload; that failure is retained. The successful arXiv v2 fallback is separately bound and is not presumed identical to its publisher version. Likewise, no equivalence between the reviewed target arXiv v3 and the 2025 version of record is asserted.\par
The current focused acceptance preserves `scientific\_decision: pending` and `human\_reviews: []`; it does not constitute human acceptance.\par
9 Y. Liu, Phase Transition of Parabolic Ginzburg–Landau Equation with Potentials of High-Dimensional Wells, arXiv:2207.12912v3, 10 May 2023, 38 pp., https://arxiv.org/abs/2207.12912v3. Frozen journal identity: Comm. Pure Appl. Math. 78(6) (2025), 1199–1247, https://doi.org/10.1002/cpa.22242. The journal and preprint are not presumed equivalent.\par
F.-H. Lin, X.-B. Pan and C.-Y. Wang, Phase Transition for Potentials of High-Dimensional Wells, Comm. Pure Appl. Math. 65(6) (2012), 833–888. Checked source: 56-page author-hosted publisher-formatted preprint, https://www.math.purdue.edu/\textasciitilde{}wang2482/lpw2.pdf, especially Section 2 and Appendix A.\par
Y. Chen and M. Struwe, Existence and Partial Regularity Results for the Heat Flow for Harmonic Maps, Math. Z. 201 (1989), 83–103. Checked primary scan: https://people.clas.ufl.edu/yun/files/paper-23.pdf.\par
G. Alberti and C. Mantegazza, A Note on the Theory of SBV Functions, Boll. Un. Mat. Ital. B (7) 11 (1997), 375–382. Checked author-hosted text: https://pagine.dm.unipi.it/alberti/ricerca/1992-96/sbv.pdf.\par
J. Fischer, T. Laux and T. M. Simon, Convergence Rates of the Allen–Cahn Equation to Mean Curvature Flow: A Short Proof Based on Relative Entropies, arXiv:2002.11994v2, 18 August 2020, https://arxiv.org/abs/2002.11994v2. The exact checked preprint has 12 pages.\par
T. Laux and Y. Liu, Nematic-Isotropic Phase Transition in Liquid Crystals: A Variational Derivation of Effective Geometric Motions, arXiv:2010.07786v2, 24 March 2021, https://arxiv.org/abs/2010.07786v2. The exact checked fallback preprint has 22 pages; its publisher version is a distinct source.\par
NYU Shanghai, Yuning Liu faculty page, selected publications, https://shanghai.nyu.edu/academics/faculty/directory/yuning-liu. Used only for public bibliographic identity and the bounded supplementary-material lookup.\par
\end{document}
